\chapter{The KES Map: Formal Specification and Convergence Conditions}
\label{app:kes}

\noindent The Kinetic-Event Synthesis (\KES) map formalizes the transition from
a current world-state to a successor hypothesis space as the primitive cognitive
operation. This appendix develops the formal specification, convergence conditions,
and relationship to the broader admissibility framework.

\section{World-States and Hypothesis Spaces}

\begin{definition}{World-State}
A \emph{world-state} $\Omega_t$ is a tuple
\[
  \Omega_t = (\mathcal{E}_t,\, \mathcal{C}_t,\, \mathcal{B}_t,\, \mathcal{H}_t)
\]
where $\mathcal{E}_t$ is the set of committed events (actualized through Pop or
Collapse operations), $\mathcal{C}_t$ is the current constraint set, $\mathcal{B}_t$
is the set of active Bind dependencies, and $\mathcal{H}_t$ is the residual
hypothesis space: the set of configurations still admissible given the committed
events and active constraints.
\end{definition}

\begin{definition}{KES Map}
The \emph{Kinetic-Event Synthesis map} is the function
\[
  \KES: \Omega_t \longrightarrow H_{t+1}
\]
where $H_{t+1}$ is the \emph{successor hypothesis space}: the set of all
world-states $\Omega_{t+1}$ that are admissible continuations of $\Omega_t$
given the constraint set $\mathcal{C}_{t+1} = \mathcal{C}_t \cup \Delta\mathcal{C}_t$,
where $\Delta\mathcal{C}_t$ represents new constraints generated by events in
$\mathcal{E}_t$.
\end{definition}

\section{The KES Algorithm}

The map $\KES(\Omega_t)$ is computed through the following sequence of operations:

\textbf{Step 1: Event propagation.} For each committed event $e \in \mathcal{E}_t$,
evaluate the constraint implications $\delta\mathcal{C}(e)$:
\[
  \Delta\mathcal{C}_t = \bigcup_{e \in \mathcal{E}_t} \delta\mathcal{C}(e).
\]

\textbf{Step 2: Constraint accumulation.} Update the constraint set:
\[
  \mathcal{C}_{t+1} = \mathcal{C}_t \cup \Delta\mathcal{C}_t.
\]
Note that constraint accumulation is monotone: $\mathcal{C}_{t+1} \supseteq
\mathcal{C}_t$, implementing the irreversibility condition.

\textbf{Step 3: Hypothesis filtration.} Generate $H_{t+1}$ by filtering all
logically possible successor states through $\mathcal{C}_{t+1}$:
\[
  H_{t+1} = \bigl\{\Omega \mid \Omega \text{ satisfies all } c \in \mathcal{C}_{t+1}\bigr\}.
\]

\textbf{Step 4: Bind propagation.} For each dependency $(e_i \to e_j) \in
\mathcal{B}_t$, verify that $H_{t+1}$ contains only states in which the
dependency is respected. Remove states that violate active bindings.

\textbf{Step 5: Residue annotation.} Annotate $H_{t+1}$ with the structural
residue $\mathcal{R}_{t+1}$: the traces of prior constraint configurations
that influenced the current filtration. This implements provenance at the
level of the hypothesis space itself.

\section{Convergence Conditions}

\begin{theorem}{KES Convergence}
Let $\{\Omega_t\}_{t \geq 0}$ be a \KES trajectory initialized at $\Omega_0$.
The trajectory \emph{converges} to a stable configuration $\Omega^*$ if and
only if there exists $T < \infty$ such that for all $t > T$:
\[
  |H_{t+1}| = |H_t|,
\]
that is, the hypothesis space stabilizes in cardinality. A necessary condition
for convergence is that the constraint accumulation rate $|\Delta\mathcal{C}_t|$
eventually vanishes: no new constraints are generated by the committed events.
\end{theorem}

\begin{proposition}{Monotone Contraction}
Under the \KES dynamics, the hypothesis space is monotonically contracting:
\[
  H_{t+1} \subseteq H_t \quad \text{for all } t.
\]
This follows directly from the monotone accumulation of constraints:
$\mathcal{C}_{t+1} \supseteq \mathcal{C}_t$ implies
$\mathcal{A}(\mathcal{C}_{t+1}) \subseteq \mathcal{A}(\mathcal{C}_t)$.
\end{proposition}

The monotone contraction property is the formal statement of why \KES implements
directed behavior without teleology: the hypothesis space narrows under constraint
accumulation, producing apparent goal-directedness as a structural consequence of
irreversibility rather than as evidence of pre-inscribed targets.

\section{Relationship to Spherepop}

The \KES map and the Spherepop calculus (Appendix~\ref{app:spherepop}) are
complementary formalizations of the same underlying process. The \KES map
describes the hypothesis space at the level of world-states; the Spherepop
operators describe the individual events that drive the \KES dynamics:

\begin{itemize}
  \item $\mathrm{Pop}(e)$ extracts an event from $H_t$ into $\mathcal{E}_t$,
    triggering constraint propagation and the next \KES step.
  \item $\mathrm{Refuse}(e)$ removes an event from $H_t$, contracting the
    hypothesis space without committing to an alternative.
  \item $\mathrm{Bind}(e_i, e_j)$ adds a dependency to $\mathcal{B}_t$,
    constraining all future filtrations.
  \item $\mathrm{Collapse}(H_t)$ commits the entire current hypothesis space
    to a definite configuration, terminating the \KES trajectory at that point.
\end{itemize}

The \KES map therefore provides the macro-level description; Spherepop provides
the micro-level mechanism.

\section{Relationship to the RSVP Fields}

Under the semantic interpretation, the \KES map corresponds to a single
time-step in the \RSVP field dynamics:
\[
  \KES(\Omega_t) \leftrightarrow
  (\Phi_{t+1}, \mathbf{v}_{t+1}, S_{t+1})
  = \text{Euler step of } \mathcal{L}_{\rm RSVP} \text{ from }
  (\Phi_t, \mathbf{v}_t, S_t).
\]
The scalar field $\Phi$ tracks the density of committed hypotheses. The vector
field $\mathbf{v}$ tracks the direction of hypothesis filtration. The entropy
field $S$ tracks the unresolved residue of the hypothesis space. Convergence
in the \KES sense corresponds to entropy relaxation in the \RSVP sense: the
system reaches a stable configuration when $S \to 0$ locally.
